
\documentclass[11pt]{article}
\usepackage[margin=1in]{geometry}
\usepackage{amsmath, amssymb, amsthm}
\usepackage{hyperref}

\title{The Parasite Machine v2:\\
Asymmetry, Representation, and Conic Logic}
\author{}
\date{}

\newtheorem{definition}{Definition}
\newtheorem{proposition}{Proposition}

\begin{document}
\maketitle

\section*{Abstract}
This document presents a refined formalization of the Parasite machine.
We construct a system in which:
\begin{itemize}
\item relations are intrinsically asymmetric and intensive,
\item counter-relations are implied rather than assumed,
\item representation functions as a curvature of implication,
\item the tetralemma becomes a phase space of relational regimes,
\item a third term (the parasite) modulates all of the above.
\end{itemize}

\section{1. From Relation to Asymmetry}
We begin by rejecting neutral relations.

\begin{definition}[Asymmetric Relation]
A primitive relation is:
\[
A \xrightarrow[\iota,\delta]{} B
\]
where:
\begin{itemize}
\item $\iota$ = intensity
\item $\delta$ = dynamic asymmetry
\end{itemize}
\end{definition}

\noindent
This implies:
\[
R(A,B) = (\iota_{AB}, \delta_{AB}) \neq (\iota_{BA}, \delta_{BA})
\]

\paragraph{Narrative}
A relation is not merely directional—it is biased, weighted, and dynamic.
Symmetry is not broken; it is never assumed.

\section{2. Implied Counter-Relation}
\begin{definition}[Mode-dependent implication]
\[
\Phi_m : R(A,B) \mapsto R(B,A)
\]
\end{definition}

\paragraph{Narrative}
The reverse relation is not given. It is produced.
This production depends on a regime $m$.

\section{3. Representation as Curvature}

\begin{definition}[Modes]
\[
\mathcal{M} = \{Id, Res, An, Con\}
\]
\end{definition}

Each mode determines implication:

\subsection*{Identity (Ellipse)}
\[
\Phi_{Id}(R(A,B)) \approx R(B,A)
\]

\paragraph{Narrative}
Strong reciprocity. Closure. Mutual implication.

\subsection*{Resemblance (Parabola)}
\[
\Phi_{Res}(R(A,B)) = \lambda R(B,A), \quad 0<\lambda<1
\]

\paragraph{Narrative}
Partial return. One-sided drift. Threshold behavior.

\subsection*{Analogy (Transformation)}
\[
\Phi_{An}(R(A,B)) = T(R(B,A))
\]

\paragraph{Narrative}
Return is transformed, not mirrored.
Structure preserved, content altered.

\subsection*{Contrariety (Hyperbola)}
\[
\Phi_{Con}(R(A,B)) = R(B,A)
\]

\paragraph{Narrative}
Return amplifies divergence.
The system splits rather than closes.

\section{4. Conic Mapping}
\[
\kappa(m) \in \{ellipse, parabola, hyperbola, mixed\}
\]

\paragraph{Narrative}
Representation is not epistemic—it is geometric.
Each mode defines curvature in relational space.

\section{5. The Parasite}
\begin{definition}[Parasite Operator]
\[
C \triangleright_m (A \xrightarrow[\iota,\delta]{} B)
\]
\end{definition}

\paragraph{Narrative}
The parasite acts on:
\begin{itemize}
\item intensity
\item asymmetry
\item implied relation
\item representational mode
\end{itemize}

It is the displaced “middle.”

\section{6. Tetralemma as Phase Space}
\[
\tau(R) \in \{K_1, K_2, K_3, K_4\}
\]

\begin{itemize}
\item $K_1$: affirmed
\item $K_2$: negated
\item $K_3$: both
\item $K_4$: neither
\end{itemize}

\paragraph{Narrative}
These are not truth values.
They are regimes of relational stabilization.

\section{7. Full System}
\begin{definition}[Parasite Machine]
\[
P = (A,B,C,R,\Phi,m,\tau,d)
\]
\end{definition}

\paragraph{Narrative}
This machine evolves:
\[
P' = \Psi(P)
\]

where $\Psi$ updates:
\begin{itemize}
\item relation $R$
\item implied relation
\item mode $m$
\item state $\tau$
\item distance $d$
\item positions $A,B,C$
\end{itemize}

\section{8. Conceptual Consequences}
\begin{itemize}
\item Asymmetry is primitive
\item Reciprocity is derived
\item Representation is curvature
\item The excluded middle becomes operative
\item Structure emerges dynamically
\end{itemize}

\section*{Final Statement}
A world can be generated from asymmetric relations, their implied returns, and a third term that modulates their curvature, without presupposing substance or conservation.

\end{document}
